\subsubsection*{Solution}
\paragraph*{Step-by-Step Derivation}

Work in natural units, $\hbar=c=1$, and assume:

\begin{itemize}[leftmargin=*]
\item canonical local density $\rho_{\rm DM}=0.4\,{\rm GeV\,cm^{-3}}$;
\item LIGO arm length $L=4\,{\rm km}$;
\item optimal projection of the vector force along the doped arm;
\item fully coherent matched-filter integration for $T=13$ years.
\end{itemize}

The density and force conventions follow the LIGO dark-photon analysis, which gives $A_0=\sqrt{2\rho_{\rm DM}}/m_A$ and $a=\epsilon e(Q_D/M)\dot A$ \href{\detokenize{https://arxiv.org/abs/1905.04316}}{Guo et al.}. LIGO's arms are each $4\,{\rm km}$ long \href{\detokenize{https://www.ligo.caltech.edu/page/ligo-gw-interferometer}}{LIGO Laboratory}.

\subparagraph*{1. Dark electric-field amplitude}

Write the nonrelativistic vector field as

\[
\mathbf A(t)=\mathbf A_0\cos(m_A t).
\]

Its energy density is

\[
\rho_{\rm DM}=\frac12 m_A^2 A_0^2,
\]

so

\[
A_0=\frac{\sqrt{2\rho_{\rm DM}}}{m_A}.
\]

The oscillating dark electric field has amplitude

\[
E_{D,0}=|\dot{\mathbf A}|_0=m_AA_0 =\sqrt{2\rho_{\rm DM}}.
\]

\subparagraph*{2. Differential acceleration}

The inner and doped outer mirrors have

\[
\left(\frac{Q_D}{M}\right)_{\rm inner} =\frac{0.5}{m_n},
\]

\[
\left(\frac{Q_D}{M}\right)_{\rm outer} =\frac{0.5+\delta q}{m_n}.
\]

Therefore the common $0.5/m_n$ contribution cancels, leaving

\[
\Delta\!\left(\frac{Q_D}{M}\right)=\frac{\delta q}{m_n}.
\]

Thus the differential acceleration amplitude is

\[
\Delta a_0 =\epsilon_{B-L}e\,\frac{\delta q}{m_n} \sqrt{2\rho_{\rm DM}}.
\]

\subparagraph*{3. Displacement and strain}

For a freely suspended mirror driven at angular frequency $\omega=m_A$,

\[
\Delta x_0=\frac{\Delta a_0}{m_A^2}.
\]

The resulting strain is

\[
h_0=\frac{\Delta x_0}{L} =\epsilon_{B-L}\delta q\, \frac{e\sqrt{2\rho_{\rm DM}}} {m_n m_A^2L}.
\]

At $f=250\,{\rm Hz}$,

\[
m_A=hf =(4.135667696\times10^{-24}\,{\rm GeV\,s})(250\,{\rm s^{-1}}) =1.033916924\times10^{-21}\,{\rm GeV}.
\]

Using

\[
1\,{\rm cm}=5.06773071616\times10^{13}\,{\rm GeV^{-1}},
\]

we obtain

\[
\rho_{\rm DM} =\frac{0.4}{(5.06773071616\times10^{13})^3} =3.07340223055\times10^{-42}\,{\rm GeV^4}.
\]

Also,

\[
m_n=0.93956542194\,{\rm GeV}, \qquad e=\sqrt{4\pi\alpha}=0.302822120769,
\]

with the neutron mass and fundamental constants taken from \href{\detokenize{https://physics.nist.gov/cuu/pdf/all.pdf}}{NIST CODATA}, and

\[
L=4000\,{\rm m} =2.02709228646\times10^{19}\,{\rm GeV^{-1}}.
\]

Consequently,

\[
\frac{e\sqrt{2\rho_{\rm DM}}}{m_nm_A^2L} =36.8757159102,
\]

so the signal strain is

\[
h_0=36.8757159102\,\epsilon_{B-L}\delta q.
\]

\subparagraph*{4. Coherent strain threshold}

For a monochromatic signal and one-sided strain noise amplitude spectral density $S_h^{1/2}$,

\[
{\rm SNR}^2=\frac{h_0^2T}{S_h}.
\]

At ${\rm SNR}=1$,

\[
h_{0,\min}=\frac{S_h^{1/2}}{\sqrt T}.
\]

The observation time is

\[
T=13(365.25)(86400) =410248800\,{\rm s}.
\]

Therefore,

\[
h_{0,\min} =\frac{3\times10^{-24}}{\sqrt{410248800}} =1.48114506204\times10^{-28}.
\]

Solving for the minimum coupling,

\[
\epsilon_{B-L,\min} =\frac{h_{0,\min}} {36.8757159102\,\delta q}.
\]

This gives

\[
\begin{array}{c|c}
\delta q & \epsilon_{B-L,\min}\\
\hline
0.074 & 5.42781972291\times10^{-29}\\
6\times10^{-3} & 6.69431099159\times10^{-28}\\
5\times10^{-4} & 8.03317318991\times10^{-27}
\end{array}
\]

The fully coherent 13-year assumption is idealized: virialized vector dark matter near $250\,{\rm Hz}$ has a coherence time of only approximately $10^3$-$10^4\,{\rm s}$, as discussed in the LIGO analysis. A realistic phase-incoherent search would therefore have a weaker, analysis-dependent reach.
